% results.tex — PRELIMINARY results (Pythia-160m, CPU). Every number traces to findings.md
% and a seeded script. Tables are structured so GPU-scaled rows drop in without restructuring.
\section{Results}
\label{sec:results}

\textbf{All results in this section are preliminary, obtained on Pythia-$160$M on CPU with $N=300$
ground-truth Pile members (seed $0$); larger-model rows are left for the GPU replication.} Every
number is reproducible from a seeded script and recorded in our results ledger.

\subsection{Membership separation is at chance on a confound-clean split}
\label{sec:res-membership}
We first reproduce, as a control, the known weakness of membership inference on pre-trained
LLMs~\cite{duan2024mia}. On a confound-clean split (members = Pile train, non-members = Pile
validation, stratified across $22$ Pile subsets to match domain), all four detectors sit at chance
at $160$M (Table~\ref{tab:membership}); on the temporally-confounded WikiMIA split the same model
shows a spurious $0.52$--$0.56$, and a $1.4$B model rises further---evidence that the WikiMIA signal
is substantially distribution shift, not membership.

\begin{table}[t]
\centering\footnotesize\setlength{\tabcolsep}{4pt}
\begin{tabular}{@{}lcccc@{}}
\toprule
Construction (model) & LOSS & Min-K\% & Min-K\%++ & zlib \\
\midrule
Pile train-vs-val, clean ($160$M) & 0.454 & 0.470 & 0.490 & 0.484 \\
WikiMIA-64, confounded ($160$M) & 0.523 & 0.539 & 0.545 & 0.564 \\
WikiMIA-64, confounded ($1.4$B) & 0.571 & 0.580 & 0.547 & 0.616 \\
\bottomrule
\end{tabular}
\caption{Membership AUC. Chance ($\approx0.5$) on the confound-clean split at $160$M; the WikiMIA
``signal'' is largely temporal/topical distribution shift. CIs in the ledger; deduplicated Pythia
gives the same chance-level result.}
\label{tab:membership}
\end{table}

\subsection{Contamination predicts leakage --- but only through loss}
\label{sec:res-headline}
Our headline analysis correlates each per-item detector score with the per-item extraction outcome
(prefix-continuation extractable memorization under greedy decoding~\cite{carlini2023quantifying}),
then controls for raw loss. Table~\ref{tab:headline} reports, for each calibrated detector, the
zero-order Spearman $\rho$, the linear partial $\rho$ given loss, the non-linear (cubic-residual)
partial $\rho$ with bootstrap CI, the FDR-corrected permutation $q$, and the mediation decomposition.

\begin{table}[t]
\centering\footnotesize\setlength{\tabcolsep}{3.5pt}
\begin{tabular}{@{}lccccc@{}}
\toprule
Detector & zero-order & partial$\mid$loss & cubic-resid.\ [95\% CI] & BH-$q$ & mediation: direct $\mid$ indirect \\
\midrule
LOSS & $+0.275$ & --- & --- & --- & (mediator) \\
Min-K\% & $+0.173$ & $-0.178$ & $-0.110$ $[-0.234,-0.002]$ & $0.058$ & $-0.394 \mid +0.567$ \\
Min-K\%++ & $+0.108$ & $-0.148$ & $-0.160$ $[-0.287,-0.041]$ & $\mathbf{0.015}$ & $-0.213 \mid +0.321$ \\
zlib & $+0.177$ & $-0.042$ & $-0.052$ $[-0.165,+0.068]$ & $0.331$ & $-0.061 \mid +0.238$ \\
\bottomrule
\end{tabular}
\caption{Headline: per-item contamination score vs.\ extraction (Spearman $\rho$), Pythia-$160$M,
$N=300$ members. The positive zero-order correlations collapse to $\approx 0$ or significantly
\emph{negative} once loss is controlled---linearly, and under the non-linear cubic-residual control
(no positive signal revives; deciles and the deduplicated arm agree). Mediation: the loss-mediated
\emph{indirect} effect is significantly positive for all three detectors while the \emph{direct}
effect is null (zlib) or negative (Min-K\%, Min-K\%++). We read this as a \emph{descriptive}
decomposition, not a causal mediation claim (see below): no calibrated detector adds positive signal
beyond loss.}
\label{tab:headline}
\end{table}

\paragraph{Collinearity caveat (why we do not over-read the negative partials).} The calibrated
detectors are deterministic transforms of the same per-token log-probabilities as loss, and are
empirically collinear with it: Spearman $\rho(\text{loss},\cdot)=0.90$ (Min-K\%), $0.74$ (Min-K\%++),
$0.74$ (zlib), with variance-inflation factors $6.2$, $2.6$, $2.4$. The strongest negative partial
(Min-K\%, the most loss-collinear detector at VIF $6.2$) is therefore consistent with a
\emph{suppression artifact} of near-collinearity rather than substantive inverse prediction; we do
not claim the calibrated detectors \emph{negatively} predict leakage. The defensible, conservative
statement is that they carry \emph{no positive} leakage signal independent of loss. Min-K\%++ and
zlib have only moderate collinearity (VIF $<3$), so their null/near-null residuals are less
attributable to collinearity.

\noindent The pre-registered decision rule asked whether any calibrated detector predicts leakage
\emph{beyond} loss (a positive partial $\rho$, CI excluding zero, FDR-significant). None does, under
the linear or the non-linear control. \textbf{Power note:} with $N=300$ and a near-degenerate
outcome ($3/300$ fully extracted), this is evidence of \emph{no positive independent signal of
appreciable size}, not proof of an exact null; the analysis is well-powered only for moderate-to-large
positive residuals, and a small positive effect at scale is not excluded (hence the GPU replication).
The per-domain breakdown (ledger) shows the loss$\leftrightarrow$extraction link is heterogeneous and
sign-flipping across domains---strongest in templated/structured domains (GitHub, StackExchange),
reversed in some prose domains (PubMed Abstracts)---so the pooled $\rho$ is a domain-mixture, not a
uniform effect.

\subsection{Extraction and PII at this scale}
\label{sec:res-extraction}
Extractable memorization is rare at $160$M: $3/300$ members are fully extractable (exact-match
extraction rate $0.010$; mean fractional extraction $0.037$), the fully-extracted items being
templated boilerplate. On the Enron-Emails-in-Pile subset we measured \emph{zero} verbatim PII
leakage ($8/36$ documents contained PII in the held suffix; none were regurgitated). We report the
PII result as a null at this scale and make no PII-exposure claim; both quantities are expected to
grow with model scale.

\subsection{Benchmark contamination (model-free $n$-gram + permutation test)}
\label{sec:res-matrix}
We complement the per-item analysis with two benchmark-level contamination tests
(Table~\ref{tab:matrix}). The model-free $n$-gram overlap against a public \emph{sample} of the Pile
($10$k documents) is a scale-invariant method but, with a sampled reference, yields only a loose
\emph{lower bound}: overlap is near-zero for MMLU ($0.2\%$ at $13$-grams), GSM8K ($0\%$), and
HumanEval ($0\%$ at $13$-grams), which certifies overlap is \emph{at least} this small and is
uninformative about true contamination---a full-Pile index (infrastructure-, not GPU-, gated) is
required for a real rate. The Oren permutation/exchangeability test~\cite{oren2024proving} at $160$M
finds the canonical ordering favoured beyond chance for MMLU ($p=0.001$) and GSM8K ($p=0.013$) but
not HumanEval ($p=0.875$); we draw \emph{no} contamination conclusion from this, as the test is
membership-based, run at sanity scale (small $k$, smallest model), and subject to a fluency/orientation
artifact---it is flagged GPU-gated and requires a fluency-control baseline before any claim.

\begin{table}[t]
\centering\footnotesize\setlength{\tabcolsep}{4pt}
\begin{tabular}{@{}lccc@{}}
\toprule
Benchmark & $13$-gram overlap (lower bound) & $8$-gram overlap & Oren $p$ ($160$M, sanity) \\
\midrule
MMLU & $0.2\%$ & $0.8\%$ & $0.001$ \\
GSM8K & $0.0\%$ & $0.0\%$ & $0.013$ \\
HumanEval & $0.0\%$ & $1.8\%$ & $0.875$ \\
\bottomrule
\end{tabular}
\caption{Benchmark-level contamination at small scale. $n$-gram cells are a \emph{lower bound}
against a $10$k Pile sample (method scale-invariant, reference under-powered); Oren $p$-values are
sanity-scale at $160$M and GPU-gated (no contamination conclusion drawn). See
\texttt{docs/contamination\_matrix.md}.}
\label{tab:matrix}
\end{table}
